\lecture{1}{26. Januar 2026}{Introduction}

\section{Introduction}

This collection of notes are the personal notes of me, Noah Hansen, for the course ``Theory of Elasticity'' (Da: Elasticitetsteori) taught on the 4th semester of the bachelor on the mechanical engineering programme at Aarhus Universitet.

The notes will take base in the lectures by Assistant Professor Matteo Pezzulla and ``Introduction to Linear Elasticity'' by Phillip L. Gould and Yuan Feng.

\subsection{Coordinate Rotation}

Given a position vector $\Vec{r}$ to a point $P$, we can resolve it into components with respect to two different Cartesian systems, $x_i$ and $x_i'$, having a common origin.

The point $P$ has coordinates $P(x_1, x_2, x_3) = P(x_i)$ in the unprimed system and $P(x_1', x_2', x_3') = P(x_i')$ in the primed system.
The linear transformation between the coordinates is then:
\begin{align*}
  x_1' &= \alpha_{11} x_1 + \alpha_{12} x_2 + \alpha_{13} x_3 \\
  x_2' &= \alpha_{21} x_2 + \alpha_{22} x_2 + \alpha_{23} x_3 \\
  x_3' &= \alpha_{31} x_1 + \alpha_{32} x_2 + \alpha_{33} x_3
\end{align*}
or
\begin{equation*}
  x_i' = \alpha_{ij} x_j
.\end{equation*}
Each of the nine quantities $\alpha_{ij}$ is the cosine of the angle between the $i$th primed and the $jth$ unprimed axis, i.e.
\begin{equation*}
  \alpha_{ij} = \cos \left( x_i', x_j  \right) = \frac{\partial x_j}{\partial x_i'} = \Vec{e}_i' \cdot \Vec{e}_j = \cos \left( \Vec{e}_i', \Vec{e}_j \right)
.\end{equation*}
These $\alpha_{ij}$'s are known as direction cosines.


\subsection{Cartesian Tensors}

A tensor of order $n$ is a set of $3^{n}$ quantities which transform from one coordinate system $x_i$ to another $x_i'$, by the following law:

\begin{table}[ht]
\centering
\caption{Tranformation laws for Tensors}
\begin{tabular}{c|c|c}
$n$ & Order & Transformation law \\
\hline
0 & Zero (scalar) & $A_i' = A_i$ \\
1 & One (vector) & $A_i' = \alpha_{ij} A_j$ \\
2 & Two (dyadic) & $A_{ij}' = \alpha_{ik}, \alpha_{jl}, A_{kl}$ \\
3 & Three & $A_{ijk}' = \alpha_{il}\alpha_{jm}\alpha_{kn} A_{lmn}$ \\
4 & Four & $A'_{ijkl} = \alpha_{im} \alpha_{jn} \alpha_{kp} \alpha_{lq} A_{mnpq}$
\end{tabular}
\end{table}

Order zero and order one tensors (scalars and vectors are already familiar quantities, wheres higher-order tensors are not.
These tensors are, however, useful to describe physical quantities with a corresponding number of associated directions.
Second-order tensors (dyadics) are especially prevalent for elasticity and the corresponding transformation may be carried out in matrix format as:
\begin{equation*}
  \Vec{A}' = \Vec{R} \Vec{A} \Vec{R}^{T}
.\end{equation*}
in which
\begin{equation*}
  \Vec{R} = \begin{bmatrix}
    \alpha_{11} & \alpha_{12} & \alpha_{13}\\
  \alpha_{21} & \alpha_{22} & \alpha_{23}\\
  \alpha_{31} & \alpha_{32} & \alpha_{33}\\
  \end{bmatrix} \qquad  \text{and} \qquad \Vec{A}' = \begin{bmatrix}
    A'_{11} & A'_{12} & A'_{13}\\
  A'_{21} & A_{22}' & A_{23}'\\
  A_{31}' & A_{32}' & A_{33}'\\
  \end{bmatrix}
.\end{equation*}


\subsubsection{Algebra of Cartesian Tensors}

Arithmetic and algebra for tensors is similar to matrix operations when it comes to addition, subtraction, equality, and scalar multiplication.
Multiplication of two tensors of order $n$ and $m$ produces a now tensor of order $n + m$, e.g.
\begin{equation*}
  A_i B_{jk} = C_{ijk}
.\end{equation*}


\subsection{Operational Tensors}

Additional tensor operations are made possible by the Kronecker delta $\delta_{ij}$, defined as,
\begin{equation*}
  \delta_{ij} \equiv \begin{cases}
    1 & i = j\\
  0 & i \neq j
  .\end{cases}
\end{equation*}
and the permutation tensor $\epsilon_{ijk}$ defined such that:
\begin{align*}
  \epsilon_{ijk} &= \frac{1}{2} (i-j)(j-k)(k-i) \\
                 &= \begin{cases}
    0 & \text{If any two of } i, j, k \text{ are equal}\\
  1 & \text{For an even permutation (forward on the number line } 1,2,3,1,2,3,\ldots ) \\
  -1 & \text{For an odd permutation (backward on the number line } 3,2,1,3,2,1,\ldots )
  .\end{cases}
\end{align*}

The Kronecker delta $\delta_{ij}$ changes the subscripts in a tensor by multiplication:
\begin{align*}
  \delta_{ij} A_i &= \delta_{1j} A_1 + \delta_{2j} A_2 + \delta_{3j} A_3 = A_j \\
  \delta_{ij} C_{ijk} &= \delta_{11} C_{11k} + \delta_{22} C_{22k} + \delta_{33} C_{33k} = C_{iik} = A_K
.\end{align*}
It is also useful in vector algebra and for coordinate transformations. E.g. starting with the dot product of two unit vectors:
\begin{equation}
  \Vec{e}_i \cdot \Vec{e}_j = \delta_{ij}
\end{equation}
and seeking the component of vector $\Vec{A}$ in the $j$-direction $A_j$ we do:
\begin{align*}
  \Vec{e}_j \cdot \Vec{A} &= \left( \Vec{e}_j \cdot \Vec{e}_i \right)A_i \\
  &= \delta_{ji} A_i \\
  &= A_j
.\end{align*}

The permutation tensor $\epsilon_{ijk}$ is useful for vector cross-product operations:
\begin{equation*}
  \epsilon_{ijk} A_j B_k \Vec{e}_i = \Vec{A} \times \Vec{B}
.\end{equation*}
